\documentclass[12pt, a4paper]{article}

\title{
    SINOPSIS\\
    MAGANG INDUSTRI SEBAGAI \textit{FRONT-END DEVELOPER}\\
    DI PT INTI UTAMA SOLUSINDO
}
\author{535230139 ALDIAN YOHANES}

% Format Font Calibri
\usepackage{fontspec}
\setmainfont{Calibri}

% Format Page Layout / Margin
\usepackage[
    a4paper, 
    left=4cm, 
    right=3cm, 
    top=4cm, 
    bottom=3cm,
]{geometry}
\pagenumbering{gobble} % Hilangkan nomor halaman


% Disable break word pakai "-"
\hyphenpenalty=10000
\exhyphenpenalty=10000
\tolerance=9999
\emergencystretch=10pt


% Default spacing ke 1.5
\usepackage{setspace}
\onehalfspacing 
\setlength{\parindent}{1.27cm}
\setlength{\parskip}{12pt}

% Import package lainnya (gambar, tabel, daftar isi)
\usepackage[indonesian]{babel}

\begin{document}
\makeatletter
\begin{center}
    \large\textbf{\@title}
\end{center}

\makeatother
PT Inti Utama Solusindo (IUS) merupakan perusahaan yang bergerak di 
bidang teknologi informasi dan menjadi bagian dari Pharos Group, 
sebuah grup perusahaan yang beroperasi dalam industri farmasi dan 
layanan kesehatan di Indonesia. Sejalan dengan transformasi digital 
pada sektor healthcare, PT IUS berperan sebagai pengembang sistem dan 
aplikasi digital yang mendukung operasional internal grup maupun 
layanan eksternal kepada fasilitas kesehatan.

Perusahaan ini mengembangkan berbagai \textit{platform} berbasis 
\textit{web} dan \textit{mobile}, seperti sistem informasi klinik 
(HealthyOne), \textit{platform} distribusi farmasi (Hospinet), 
serta aplikasi pendukung operasional pemasaran dan \textit{marketplace} 
internal (Canvasser). Fokus utama kegiatan perusahaan adalah 
pengembangan dan pemeliharaan sistem informasi terintegrasi 
yang mendukung efisiensi, akurasi data, dan optimalisasi layanan 
kesehatan berbasis teknologi.

Kegiatan magang industri dilaksanakan di kantor PT Inti Utama Solusindo 
yang berlokasi di Gedung Century, Jl. Limo No.6/45, RT.14/RW.02, 
Grogol Utara, Kec. Kebayoran Lama, Kota Jakarta Selatan. 
Program magang berlangsung selama satu tahun, dimulai pada 
12 Januari 2026 hingga 12 Januari 2027.

Pelaksanaan kerja dilakukan secara penuh di kantor (\textit{Work From 
Office}/WFO) dengan jadwal kerja Senin sampai Jumat pukul 08.00–17.00 
WIB. Skema magang ini termasuk dalam program MBKM Merdeka 
Belajar Kampus Merdeka Fakultas Teknologi Informasi Universitas 
Tarumanagara dalam semester ke-5 yang dilanjutkan dengan konversi 
mata kuliah ke magang industri pada semester ke-7.

Magang industri ini dilaksanakan sebagai bagian dari implementasi 
program MBKM untuk memberikan pengalaman profesional langsung kepada 
mahasiswa dalam lingkungan kerja nyata. Tujuan magang secara khusus 
meliputi peningkatan kompetensi teknis di bidang pengembangan 
perangkat lunak, terutama \textit{Front-End Development} menggunakan 
Flutter dan React, penerapan teori perkuliahan seperti 
\textit{Software Development Life Cycle} (SDLC), \textit{Object-Oriented 
Programming} (OOP), \textit{clean code}, dan prinsip SOLID dalam 
proyek industri nyata, pengembangan kemampuan kolaborasi 
dalam tim IT yang terdiri dari Project Manager, Quality Assurance, 
dan \textit{developer} lainnya, pemahaman \textit{project 
management system} serta \textit{version control} profesional 
menggunakan OpenProject dan GitLab, serta peningkatan kesiapan 
kerja, profesionalisme, dan kemampuan adaptasi 
terhadap standar industri teknologi informasi.

Bidang kegiatan utama PT Inti Utama Solusindo mencakup pengembangan 
\textit{sistem} dan \textit{aplikasi} berbasis teknologi informasi untuk mendukung 
\textit{ekosistem} layanan kesehatan dan farmasi. Aktivitas tersebut meliputi 
pengembangan \textit{aplikasi} \textit{sistem informasi klinik} berbasis \textit{web} dan 
\textit{mobile}, pembangunan \textit{platform} distribusi dan \textit{marketplace} farmasi digital,
pemeliharaan \textit{sistem}, perbaikan \textit{bug}, dan pengembangan \textit{fitur} baru sesuai 
kebutuhan bisnis, serta integrasi \textit{sistem} berbasis \textit{API} dan manajemen 
\textit{database}. Sebagai \textit{Front-End Developer Intern}, fokus \textit{kerja} adalah 
pada pengembangan \textit{user interface}, migrasi \textit{aplikasi} Flutter ke versi 
terbaru, serta pengembangan ulang \textit{sistem} berbasis React dengan 
pendekatan \textit{modular} dan \textit{reusable component}.

\end{document}